\documentclass[10pt,a4paper]{amsart}
\usepackage[margin=19mm]{geometry}
\usepackage{amsmath,amssymb,mathtools}
\usepackage[scr=rsfs,cal=euler]{mathalfa}
\usepackage{xfrac}
\usepackage[unicode,hidelinks,bookmarks=false]{hyperref}
\newcommand{\ie}{i.e.\ }
\newcommand{\cf}{cf.\ }
\newcommand{\resp}{respectively\ }
\newcommand{\Subsection}{\subsection}
\providecommand{\nobreakdash}{\nobreak\hbox{-}}
\setlength{\emergencystretch}{2em}
\multlinegap0pt
\usepackage{amssymb}
\usepackage{mathtools}
\usepackage{cancel}
\usepackage{xcolor}
\usepackage{algorithm}\usepackage{algpseudocode}
\usepackage{float}
\usepackage{tikz}
\newtheorem{lemma}{Lemma}
\newtheorem{theorem}{Theorem}
\newcommand{\N}{\mathbb{N}}
\newcommand{\aset}[1]{\left \langle #1 \right \rangle}
\newcommand{\bset}[1]{\left \lbrace #1 \right \rbrace}
\title{Karhunen Lo\`eve Expansions of Hilbert Space-Valued Random Elements}
\author{Trajan Murphy}
\date{Source: arXiv:2604.12042v1 (CC BY 4.0)}
\begin{document}
\maketitle

\noindent\emph{Source and licence.} Karhunen Lo\`eve Expansions of Hilbert Space-Valued Random Elements. Version arXiv:2604.12042v1 (CC BY 4.0), distributed by arXiv under the Creative Commons Attribution 4.0 International licence, http://creativecommons.org/licenses/by/4.0/, as recorded in the arXiv licence metadata for this exact version and shown on the abstract page https://arxiv.org/abs/2604.12042v1. Full licence notice: \texttt{licenses/arxiv-2604.12042-CC-BY-4.0.txt}.\par\smallskip

\noindent\textit{Editorial scope.} Counted unit: the article's theorem ONB (source0.tex lines 634--636). The article's complete original proof of it is delivered with it (source0.tex lines 638--684, one \texttt{proof} environment of 47 lines). No mathematics is written, repaired or re-derived: the statement and the proof are the article's own lines, in the article's own notation, and no retained line is substituted or re-flowed.  Retained uncounted as the source-local context the proof needs: lines 620--622.  Nothing was omitted from any retained span, so the package carries no omission record.  No retained line contains a citation, so the wrapper carries no bibliography and no bibliography entry.  No retained line contains a figure environment, an image-inclusion command or drawing code, so no image is delivered and the package declares no asset dependency.  Editorial formatting only, in addition: an AMS \texttt{amsart} preamble that restates the article's own declarations and macros used by the retained lines, namely the environments \texttt{lemma}, \texttt{theorem} on separate counters, together with the standard packages those lines need.  The retained environments print in this document as Lemma 1, Theorem 1 in order of appearance, whereas the article prints the same items with the numbering of its own full document.  The complete native source of the article as distributed in the arXiv e-print for this exact version is bundled unchanged as \texttt{sources/198378/source0.tex}.
\begin{lemma} \label{He > 0 f.o}
    If $He \neq 0$ for uncountably many $e \in \mathscr E$, then there exists an $\varepsilon > 0$ such that $\|He\|_\mathcal G > \varepsilon$ for infinitely many $e \in \mathscr E$. 
\end{lemma}

\begin{theorem} \label{He > 0 f.o. regardless of ONB}
    Suppose $\sum_{e \in \mathscr E} \|He\|_\mathcal G^2 < \infty$ for one orthonormal basis $\mathscr E$ of $\mathcal H$. Then $Hf = 0$ for all but at most countably many $f \in \mathscr F$ for any other orthonormal basis $\mathscr F$ of $\mathcal H$
\end{theorem}

\begin{proof}
    Consider another orthonormal basis $\mathscr F$ of $\mathcal H$. Suppose instead that $\|Hf\|_\mathcal G > 0$ for uncountably many $f \in \mathscr F$. By Lemma \ref{He > 0 f.o}, we may find an $\varepsilon > 0$ and a countably infinite set $\mathscr F_0$ of distinct unit vectors in $\mathscr F$ such that $\|Hf\|^2_\mathcal G > \varepsilon$ for each $f \in \mathscr F_0$.

    Define $\mathscr E_0 := \bset{e \in \mathscr E: He \neq 0}$, and let $e_k$ be an enumeration of $\mathscr E_0$. For each $k \in \N$, define $\mathscr F_k := \bset{f \in \mathscr F: \aset{f, e_k} \neq 0} $. Finally, define $\mathscr F_\infty := \bigcup_{k=0}^\infty \mathscr F_k$. As each $\mathscr F_k$ is countable, so is $\mathscr F_\infty$. 
    
   Observe that for any distinct indices, $k, k' \in \N$ that 

   $$\delta_{kk'} = \aset{e_k, e_{k'}}_\mathcal H = \sum_{f \in \mathscr F} \aset{e_k, f}_\mathcal H \aset{e_{k'}, f}_\mathcal H = \sum_{f \in \mathscr F_\infty} \aset{e_k, f}_\mathcal H \aset{e_{k'}, f}_\mathcal H$$

   which holds since, by construction, $\mathscr F_\infty$ contains all the vectors $f$ for which $\aset{e_k,f}_\mathcal H$ and $\aset{e_{k'}, f}_\mathcal H$ are nonzero. 
    
    Now observe that 

    \begin{align*}
        \sum_{e \in \mathscr E} \|He\|^2_\mathcal G
        &= 
         \sum_{e \in \mathscr E_0} \|He\|^2_\mathcal G
         \\&=
          \sum_{e, e' \in \mathscr E_0} \aset{He, He'}_{\mathcal G} \delta_{ee'}
            \\&=
          \sum_{e, e' \in \mathscr E_0} \aset{He, He'}_{\mathcal G} 
          \sum_{f \in \mathscr F_\infty} \aset{f,e}_\mathcal H \aset{f,e'}_\mathcal H
    \end{align*}

    Now let $f_k$ be an enumeration of $\mathscr F_\infty$ so that we may write 

    $$ \sum_{e \in \mathscr E} \|He\|^2_\mathcal G = \sum_{e, e' \in \mathscr E_0} \aset{He, He'}_{\mathcal G} 
          \sum_{k = 1}^\infty \aset{f_k,e}_\mathcal H \aset{f_k,e'}_\mathcal H $$

    On the other hand let $\bset{k_n}_{n=1}^\infty$ be the indices $k$ for which $\|Hf_k\|_\mathcal G > \varepsilon$. We have 

    \begin{align*}
        N\varepsilon^2 
        & \leq 
        \sum_{k_n < N} \|Hf_{k_n}\|_\mathcal G^2
        \\&\leq 
        \sum_{k < N} \|Hf_k\|_\mathcal G^2
          \\& = 
        \sum_{k < N} \left \| \sum_{e \in \mathcal E} \aset{f_k,e}_\mathcal H  He \right\|_\mathcal G^2
          \\& = 
        \sum_{k < N} \sum_{e,e' \in \mathcal E_0}  \aset{f_k,e}_\mathcal H  \aset{f_k,e'}_\mathcal H \aset{He, He'}_\mathcal G
          \\& = 
       \sum_{e,e' \in \mathcal E_0}\aset{He, He'}_\mathcal G  \sum_{k < N}  \aset{f_k,e}_\mathcal H  \aset{f_k,e'}_\mathcal H  
    \end{align*}

Taking $N \to \infty$ on the last line equals $\|He\|^2_\mathcal G$, but this is a contradiction since $N\varepsilon^2$ is unbounded.   
\end{proof}

\end{document}
